\section{Directed rays of an elementary fan}
\label{sec:area_law_cell_fan_rays}

Let $o\in\R^2$, $\ell\in\N$ and $z\in\Z^2$ be arbitrary, and let
$Q=o+2^\ell(z+[0,1]^2)$ be the corresponding closed dyadic square.
Write $c$ for its center. On each side choose either the whole side or
its two midpoint halves, and let $\mathcal E$ index the resulting
ordered elementary bases $[a_i,b_i]$.
By Theorem~\ref{thm:area_law_cell_fan_base_radius}, every final endpoint
satisfies $\|b_i-c\|_\infty=2^\ell/2>0$.

For nonzero vectors $u,w\in\R^2$, write $u\sim w$ if
$w=\lambda u$ for some $\lambda>0$. Let
$\mathcal R=(\R^2\setminus\{0\})/{\sim}$ be the set of directed rays,
and write $[u]$ for the class of $u$.
Thus opposite nonzero vectors represent distinct directed rays.
These are the directions from a cell center in the fan construction of
\cite[Section~11]{OpenAI2026AreaLaw}.

% Source: 10-geometry.tex, prop:two-families, lines 299--323,
% and geometry:initial-stars, lines 333--370;
% revision adc7f1241b42e322a6451854ab7e4b4c146bf78a.
% Manuscript: preprints/A-two-dimensional-area-law-from-a-global-spectral-gap-September-24-2026/
% build/sections/10-geometry.tex.

\begin{theorem}[Equality of endpoint directions]
    \label{thm:area_law_cell_fan_end_same_ray}
    \lean{TNLean.PEPS.AreaLaw.Geometry.cellFanEnd_sameRay_iff}
    \leanok
    \uses{def:area_law_cell_fan}
    For every $i,j\in\mathcal E$,
    \begin{align}
        b_i-c\sim b_j-c & \quad\Longleftrightarrow\quad i=j.
    \end{align}
\end{theorem}
\begin{proof}\leanok
    \uses{thm:area_law_cell_fan_base_radius,
        thm:area_law_cell_fan_final_endpoint_injective}
    The two vectors have the same positive norm $r=2^\ell/2$.
    If $b_j-c=\lambda(b_i-c)$ with $\lambda>0$, homogeneity of the
    norm gives $r=\lambda r$, hence $\lambda=1$. Thus $b_i=b_j$.
    Injectivity of final endpoints gives $i=j$.
    Conversely, equal indices give equal vectors, so the positive
    multiple $\lambda=1$ proves that they have the same direction.
\end{proof}

\begin{definition}[Directed ray of an elementary base]
    \label{def:area_law_cell_fan_ray}
    \lean{TNLean.PEPS.AreaLaw.Geometry.cellFanRay}
    \leanok
    \uses{def:area_law_cell_fan,thm:area_law_cell_fan_base_radius}
    Assign to each elementary base its final directed ray from the
    center:
    \begin{align}
        \rho(i) & =[b_i-c]\in\mathcal R.
    \end{align}
    The positive common radius makes every vector in this definition
    nonzero. Only positive multiples are identified: $[u]\ne[-u]$
    for every $u\ne0$.
\end{definition}

\begin{theorem}[Injectivity of the directed-ray assignment]
    \label{thm:area_law_cell_fan_ray_injective}
    \lean{TNLean.PEPS.AreaLaw.Geometry.cellFanRay_injective}
    \leanok
    \uses{def:area_law_cell_fan_ray}
    The map $\rho:\mathcal E\to\mathcal R$ is injective.
\end{theorem}
\begin{proof}\leanok
    \uses{def:area_law_cell_fan_ray,thm:area_law_cell_fan_end_same_ray}
    Equality of the two ray classes means that their nonzero
    representatives have the same direction. The endpoint-direction
    equivalence then gives equality of their indices.
\end{proof}
